\documentclass[10pt]{beamer}
\usetheme{metropolis}
\usepackage{graphicx}
\usepackage{listings}
\usepackage{booktabs}
\usepackage{xcolor}
\definecolor{codebg}{HTML}{F5F5F5}
\definecolor{codegreen}{HTML}{2E7D32}
\definecolor{codegray}{HTML}{757575}
\definecolor{codeblue}{HTML}{1565C0}
\lstdefinestyle{rstyle}{language=R,backgroundcolor=\color{codebg},basicstyle=\ttfamily\scriptsize,
  keywordstyle=\color{codeblue}\bfseries,stringstyle=\color{codegreen},
  commentstyle=\color{codegray}\itshape,breaklines=true,frame=single,
  rulecolor=\color{codegray!30},showstringspaces=false,tabsize=2,
  morekeywords={library,prcomp,summary,biplot,autoplot,ggfortify,Rtsne,
    umap,ggplot,aes,geom_point,geom_segment,geom_text,scale_color_manual,
    theme_minimal,labs}}
\lstdefinestyle{pythonstyle}{language=Python,backgroundcolor=\color{codebg},basicstyle=\ttfamily\scriptsize,
  keywordstyle=\color{codeblue}\bfseries,stringstyle=\color{codegreen},
  commentstyle=\color{codegray}\itshape,breaklines=true,frame=single,
  rulecolor=\color{codegray!30},showstringspaces=false,tabsize=4,
  morekeywords={import,from,as,pd,np,plt,PCA,TSNE,UMAP,StandardScaler,
    fit_transform,explained_variance_ratio_,components_}}
\title{Dimensionality Reduction Visualization}
\subtitle{Workshop 5 --- Module 3}
\author{Prof.Asoc.\ Endri Raco}
\institute{Data Visualization for Data Scientists \\ UNYT Tirana}
\date{2025/26}
\begin{document}

\begin{frame}
  \titlepage
\end{frame}

\begin{frame}{Module Roadmap}
  \begin{enumerate}
    \item Why reduce: making the invisible visible
    \item PCA: linear projection onto maximum-variance axes
    \item The scree plot: how many components to keep
    \item The biplot: scores + loadings in one figure
    \item t-SNE: nonlinear, local structure preservation
    \item UMAP: nonlinear, global + local topology
  \end{enumerate}
  \begin{exampleblock}{Learning Outcome}
    You will interpret PCA biplots (loading arrows = variable
    contribution), produce scree plots to choose the number of
    components, and compare PCA / t-SNE / UMAP projections
    to understand their different preservation properties.
  \end{exampleblock}
\end{frame}

\section{Why Reduce?}

\begin{frame}{Make the Invisible Visible}
  \begin{center}
    \includegraphics[width=0.88\textwidth]{figures_w05m03/why_reduce}
  \end{center}
  \begin{alertblock}{Pro-Tip}
    Dimension reduction is \textbf{not} about discarding information.
    It is about finding the 2D view that preserves the most important
    structure (variance for PCA, neighbourhoods for t-SNE, topology
    for UMAP). Think of it as choosing the best camera angle for a
    3D object.
  \end{alertblock}
\end{frame}

\section{PCA}

\begin{frame}{PCA: Linear Projection onto Max-Variance Axes}
  \begin{center}
    \includegraphics[width=0.72\textwidth]{figures_w05m03/pca_biplot}
  \end{center}
\end{frame}

\begin{frame}{How to Read a Biplot}
  \centering
  \begin{tabular}{lp{6cm}}
    \toprule
    \textbf{Element} & \textbf{Interpretation} \\
    \midrule
    Points (scores) & Observations projected into PC space. Close points = similar multivariate profiles. \\
    Arrows (loadings) & Original variables. Arrow direction = which PC the variable contributes to. Arrow length = strength of contribution. \\
    Angle between arrows & $\cos(\theta) \approx$ correlation. Parallel = positive corr. Opposite = negative. Perpendicular = uncorrelated. \\
    Distance from origin & Extremity. Points far from centre are unusual in PC space (potential outliers). \\
    \% variance on axes & How much of total variance this view captures. PC1+PC2 $> 70\%$ = good 2D summary. \\
    \bottomrule
  \end{tabular}
\end{frame}

\begin{frame}{Scree Plot: How Many Components?}
  \begin{center}
    \includegraphics[width=0.68\textwidth]{figures_w05m03/scree}
  \end{center}
  \begin{exampleblock}{Data Insight}
    Two rules of thumb: (1) \textbf{Elbow rule}: keep components before
    the scree ``flattens.'' (2) \textbf{80\% rule}: keep enough
    components to explain $\geq$80\% of cumulative variance. Here, 2--3
    components suffice for the 6D data.
  \end{exampleblock}
\end{frame}

\begin{frame}[fragile]{PCA in Code}
  \begin{columns}[T]
    \begin{column}{0.48\textwidth}
      \textbf{R}
\begin{lstlisting}[style=rstyle]
# PCA
pca <- prcomp(df[, num_cols],
  scale. = TRUE)  # always scale!

# Scree
summary(pca)  # proportion of var

# Quick biplot
biplot(pca)   # base R

# ggplot biplot (ggfortify)
library(ggfortify)
autoplot(pca,
  data = df,
  colour = "Cluster",
  loadings = TRUE,
  loadings.label = TRUE,
  loadings.colour = "#E65100") +
  theme_minimal()
\end{lstlisting}
    \end{column}
    \begin{column}{0.48\textwidth}
      \textbf{Python}
\begin{lstlisting}[style=pythonstyle]
from sklearn.decomposition import PCA
from sklearn.preprocessing import (
    StandardScaler)

X_scaled = StandardScaler().fit_transform(
    df[num_cols])
pca = PCA(n_components=2)
scores = pca.fit_transform(X_scaled)

# Scree
var_exp = pca.explained_variance_ratio_

# Biplot (manual)
fig, ax = plt.subplots()
ax.scatter(scores[:,0], scores[:,1],
    c=labels, cmap="viridis", s=15)
# Loading arrows
loadings = pca.components_.T * 3
for i, name in enumerate(num_cols):
    ax.annotate("", xy=loadings[i],
        xytext=(0,0), arrowprops=
        dict(arrowstyle="->",
             color="#E65100"))
    ax.text(*loadings[i]*1.1, name)
\end{lstlisting}
    \end{column}
  \end{columns}
\end{frame}

\section{t-SNE}

\begin{frame}{t-SNE: Nonlinear, Local Structure}
  \begin{center}
    \includegraphics[width=0.95\textwidth]{figures_w05m03/tsne_perplexity}
  \end{center}
  \begin{alertblock}{Pro-Tip}
    t-SNE (van der Maaten \& Hinton, 2008) preserves \textbf{local}
    neighbourhoods: nearby points stay nearby. But \textbf{distances
    between clusters are meaningless} --- cluster size and gap width
    are artefacts. The \textbf{perplexity} hyperparameter ($\sim$5--50)
    controls the effective neighbourhood size. Always try 2--3 values.
  \end{alertblock}
\end{frame}

\begin{frame}{t-SNE: What You Can and Cannot Conclude}
  \centering
  \begin{tabular}{lcc}
    \toprule
    \textbf{Conclusion} & \textbf{Valid?} & \textbf{Why} \\
    \midrule
    ``These points form a cluster'' & \textcolor{codegreen}{\checkmark} & Local structure preserved \\
    ``Cluster A is closer to B than C'' & \textcolor{red}{\texttimes} & Inter-cluster distances arbitrary \\
    ``Cluster A is larger than B'' & \textcolor{red}{\texttimes} & Size is an artefact of density \\
    ``This point is an outlier'' & \textcolor{codegreen}{\checkmark} & Isolated points = dissimilar \\
    ``PC1 correlates with variable X'' & \textcolor{red}{\texttimes} & No axes interpretation (unlike PCA) \\
    \bottomrule
  \end{tabular}
\end{frame}

\section{UMAP}

\begin{frame}{PCA vs t-SNE vs UMAP: Same Data, Three Views}
  \begin{center}
    \includegraphics[width=0.95\textwidth]{figures_w05m03/three_methods}
  \end{center}
\end{frame}

\begin{frame}[fragile]{t-SNE and UMAP in Code}
  \begin{columns}[T]
    \begin{column}{0.48\textwidth}
      \textbf{R}
\begin{lstlisting}[style=rstyle]
# t-SNE
library(Rtsne)
tsne <- Rtsne(X_scaled,
  dims = 2, perplexity = 30,
  check_duplicates = FALSE)
plot(tsne$Y, col = labels, pch = 16)

# UMAP
library(umap)
um <- umap(X_scaled)
plot(um$layout, col = labels, pch=16)

# ggplot wrapper
tibble(x = tsne$Y[,1],
       y = tsne$Y[,2],
       cluster = factor(labels)) |>
  ggplot(aes(x, y, color = cluster)) +
  geom_point(alpha = 0.5, size = 1) +
  theme_minimal() +
  labs(title = "t-SNE (perp=30)")
\end{lstlisting}
    \end{column}
    \begin{column}{0.48\textwidth}
      \textbf{Python}
\begin{lstlisting}[style=pythonstyle]
# t-SNE
from sklearn.manifold import TSNE
tsne = TSNE(n_components=2,
    perplexity=30, random_state=42)
X_tsne = tsne.fit_transform(X_scaled)
plt.scatter(X_tsne[:,0], X_tsne[:,1],
    c=labels, cmap="viridis", s=10)

# UMAP
# pip install umap-learn
from umap import UMAP
um = UMAP(n_components=2,
    random_state=42)
X_umap = um.fit_transform(X_scaled)
plt.scatter(X_umap[:,0], X_umap[:,1],
    c=labels, cmap="viridis", s=10)

# Key hyperparameters:
# t-SNE: perplexity (5-50)
# UMAP:  n_neighbors (5-50),
#        min_dist (0.0-1.0)
\end{lstlisting}
    \end{column}
  \end{columns}
\end{frame}

\section{Choosing a Method}

\begin{frame}{PCA vs t-SNE vs UMAP}
  \begin{center}
    \includegraphics[width=0.95\textwidth]{figures_w05m03/comparison}
  \end{center}
\end{frame}

\begin{frame}{Decision Rules}
  \centering
  \begin{tabular}{lll}
    \toprule
    \textbf{Goal} & \textbf{Method} & \textbf{Why} \\
    \midrule
    Understand variable contributions & PCA & Loadings are interpretable \\
    Find the ``best 2D summary'' & PCA & Max-variance projection \\
    Discover clusters visually & t-SNE or UMAP & Nonlinear separation \\
    Large dataset ($>$50K rows) & UMAP & $O(n \log n)$ vs $O(n^2)$ \\
    Need reproducibility & PCA & Deterministic \\
    Preserve global distances & PCA or UMAP & t-SNE distorts global \\
    \bottomrule
  \end{tabular}
  \vspace{6pt}
  \begin{alertblock}{Pro-Tip}
    \textbf{Best practice}: start with PCA (interpretable, fast), check
    the scree plot. If the first 2 PCs capture $<$50\% variance, the
    data has complex nonlinear structure --- then try t-SNE or UMAP.
    Always show PCA alongside t-SNE/UMAP so the reader can compare.
  \end{alertblock}
\end{frame}

\section{Common Pitfalls}

\begin{frame}{Five Pitfalls of Dimension Reduction}
  \centering
  \begin{tabular}{rlp{5cm}}
    \toprule
    \textbf{\#} & \textbf{Pitfall} & \textbf{Solution} \\
    \midrule
    1 & Forgetting to scale & Always \texttt{scale. = TRUE} (R) or \texttt{StandardScaler()} (Python) \\
    2 & Over-interpreting t-SNE distances & State: ``distances between clusters are not meaningful'' \\
    3 & Single perplexity / n\_neighbors & Try 3+ values; show in a panel \\
    4 & Treating PCs as ``real'' variables & PCs are combinations; interpret via loadings \\
    5 & Running t-SNE on $>$50K rows & Use UMAP or subsample first \\
    \bottomrule
  \end{tabular}
\end{frame}

\section{Synthesis}

\begin{frame}{Key Takeaways}
  \begin{enumerate}
    \item \textbf{PCA} = linear, interpretable (biplot + scree),
          deterministic. Use when you need to understand \emph{which
          variables drive} the structure. Always scale first.
    \item \textbf{Scree plot}: elbow rule or 80\% cumulative variance
          to choose the number of components.
    \item \textbf{Biplot}: points = observations, arrows = variable
          contributions, angle between arrows $\approx$ correlation.
    \item \textbf{t-SNE}: nonlinear, excellent at revealing clusters.
          But distances between clusters are meaningless; always
          report perplexity.
    \item \textbf{UMAP}: nonlinear like t-SNE but faster and better at
          preserving global structure. Preferred for large datasets.
    \item \textbf{Best practice}: PCA first (interpretable baseline),
          then t-SNE/UMAP if nonlinear structure suspected.
  \end{enumerate}
\end{frame}

\begin{frame}{Essential Readings}
  \begin{itemize}
    \item Jolliffe, I.~T. (2002). \emph{Principal Component Analysis},
          2nd ed. Springer.
    \item van der Maaten, L. \& Hinton, G. (2008). ``Visualizing Data
          using t-SNE.'' \emph{JMLR}, 9, 2579--2605.
    \item McInnes, L. et al. (2018). ``UMAP: Uniform Manifold
          Approximation and Projection.'' arXiv:1802.03426.
    \item Wattenberg, M. et al. (2016). ``How to Use t-SNE Effectively.''
          \emph{Distill}. \texttt{https://distill.pub/2016/misread-tsne/}
  \end{itemize}
  \vspace{6pt}
  \textbf{Next Module}: Heatmaps \& Clustered Displays (W05-M04)
\end{frame}

\end{document}
